\documentclass{situationist-preprint}

\title{The Inverted Spectrum and Functionalism}
\msid{MND-22-a}
\msfamily{MND-22}
\msdiscipline{Philosophy of Mind}
\mstype{Principal manuscript}
\mskeywords{inverted spectrum; qualia; functionalism; quality space; undetectability; Russellian monism}

\begin{document}
\maketitle

\begin{abstract}
Could two subjects be functionally alike while the colors that look red to one look green to the other? The inverted spectrum has been used both to refute functionalism and, by verificationists, to show that qualia talk is empty. We argue that inversion is possible but undetectable in principle, and that it therefore has no consequence for functionalist theories of mind. Possibility follows from three claims defended elsewhere: qualia are intrinsic, their structure is relational, and functional organization is multiply realizable in physical systems whose intrinsic natures may differ. Undetectability follows from the fact that every test, behavioral or neural, probes relations, and an inversion that preserves all relations makes no difference to any of them. Functionalism is untouched because it individuates mental states by their functional roles, and an inverted pair occupy the same roles: they see, believe, discriminate and desire alike. What the inversion shows is that qualitative character is a property of the realizers of those roles, as the intrinsic view of qualia already holds.
\end{abstract}

\section{The problem}

The possibility was raised by \citet[II.xxxii.15]{Locke1690}: the idea that a violet produces in one man's mind might be the idea a marigold produces in another's, without either being able to tell. In the functionalist era the case became an argument. If two subjects could be functionally identical while their color experiences were inverted, qualitative character is not a functional property, and functionalism about qualia is false \citep{Block1990}. \citet{Shoemaker1982} replied that qualia are functionally definable as a class even if inversion within the class cannot be ruled out. \citet{Dennett1988} argued in the other direction that, since no test could reveal an inversion, the hypothesis is empty, and so are the qualia it presupposes. \citet{Hardin1988} and \citet{Palmer1999} argued that the asymmetries of human color space constrain which inversions could go undetected, so that many candidate inversions would show up in behavior.

We argue that the three questions, whether inversion is possible, whether it could be detected, and whether it matters for functionalism, have three answers that fit together: it is possible, it is undetectable in principle, and it does not matter for functionalism as a theory of mind.

\section{Framework}

\begin{definition}\label{d:space}
A \emph{quality space} is a set of qualities together with the relations among them that a subject's discriminations track: similarity, betweenness, opponency, relative lightness and the like. A \emph{relation-preserving inversion} is a mapping of a quality space onto itself, other than the identity, that preserves all those relations.
\end{definition}

\begin{definition}\label{d:inv}
Two subjects are \emph{spectrum-inverted} relative to each other if they are functionally alike and, for some relation-preserving inversion $\iota$, whenever one has an experience with quality $q$ in a given circumstance, the other has an experience with quality $\iota(q)$ in that circumstance.
\end{definition}

\begin{assumption}[Intrinsic qualia]\label{a:intr}
Qualia are intrinsic properties of experiences; their relations to one another can be described relationally, but a relational description does not fix which quality occupies a given place in the space. This is the result of MND-10.
\end{assumption}

\begin{assumption}[Variable realization]\label{a:real}
A functional organization can be realized by physical systems whose intrinsic natures differ, and the qualities of experience are grounded in the intrinsic natures of their physical realizers. The first conjunct is the multiple realizability on which MND-08 relies; the second is the thesis of MND-01.
\end{assumption}

\begin{assumption}[Symmetry]\label{a:sym}
Some possible quality space for color admits a relation-preserving inversion.
\end{assumption}

Assumption~\ref{a:sym} is weaker than the claim that human color space admits one. Palmer's analysis shows that the asymmetries of our space rule out many candidates; whether some approximate inversion survives is an empirical question we leave open. A creature whose color space is more symmetric than ours is enough for the possibility in question.

\section{The argument}

\begin{proposition}\label{p:poss}
Spectrum inversion is possible.
\end{proposition}

\begin{proof}
Take a color space and a relation-preserving inversion $\iota$, by Assumption~\ref{a:sym}. Take a subject whose experiences have qualities in that space. By Assumption~\ref{a:real}, the functional organization of the subject can be realized by a second physical system with different intrinsic natures, and the qualities of its experiences are grounded in those natures. By Assumption~\ref{a:intr}, the organization does not fix which quality occupies each place, so nothing excludes a realizer whose natures ground $\iota(q)$ wherever the first grounds $q$. Since $\iota$ preserves all relations, the second subject's discriminations stand in the same relations as the first's, as functional likeness requires. The two are spectrum-inverted.
\end{proof}

The argument does not appeal to the conceivability of inversion. MND-05 holds that, for phenomenal concepts, ideal conceivability does not entail possibility, and the proof respects this: possibility is derived from the multiple realizability of function and the grounding of quality in intrinsic nature. Note also that the two subjects differ physically, in the intrinsic natures of their realizers; a physical duplicate with inverted qualia is not possible on MND-01, for the same reason that a zombie is not.

\begin{proposition}\label{p:undet}
No test, behavioral or neural, can detect a relation-preserving inversion between two subjects.
\end{proposition}

\begin{proof}
A behavioral test records discriminations, matches, verbal reports and reactions; each is a response to stimuli that depends on the relations among the qualities the stimuli produce, and on the functional organization through which those relations issue in behavior. A neural test records the structural and dynamical profile of the realizers, which physics and neuroscience describe relationally. MND-10 establishes that any inversion preserving all relations is undetectable by any relational test, and every test of these kinds is relational. Introspection does not help: each subject introspects only her own qualities, and her reports of them are the same as the other's.
\end{proof}

\begin{theorem}\label{t:main}
Spectrum inversion is possible, undetectable in principle, and without consequence for functionalist theories of mind.
\end{theorem}

\begin{proof}
The first two conjuncts are Propositions~\ref{p:poss} and~\ref{p:undet}. For the third: functionalism holds that mental states are individuated by their functional roles, and MND-08 that thinking is a functional property. Spectrum-inverted subjects are functionally alike by Definition~\ref{d:inv}, so they are in the same mental states under every functional individuation: both see the ripe tomato as red, believe it is red, can tell it from the leaves, and want the ripe ones. The inversion concerns a property that functionalism does not use to individuate any state: the intrinsic quality of the realizer of a role. Functionalism is therefore neither refuted nor confirmed by it.
\end{proof}

\section{Objections and replies}

\begin{objection}[Functionalism about qualia]
Some functionalists hold that qualitative character is itself functional. Proposition~\ref{p:poss} refutes them, so inversion has consequences for some functionalist theories.
\end{objection}

\begin{reply}
It has consequences for a functionalist theory of qualitative character, which MND-10 already rejects in holding qualia intrinsic. The thesis is about functionalist theories of mind: of what it is to be in a mental state, to perceive or to think. On those theories the inverted subjects have the same mind, and the difference between them is located where Shoemaker located it, in the realizers of functionally defined states.
\end{reply}

\begin{objection}[Verificationism]
A difference that no possible test could detect is no difference \citep{Dennett1988}.
\end{objection}

\begin{reply}
The principle is false in general: it would exclude the intrinsic natures of physical entities, which physics does not detect because it describes only their structure, and MND-01 argues that these natures are real. Undetectability by relational tests is what we should expect of intrinsic properties, not evidence against them.
\end{reply}

\begin{objection}[Intrasubjective inversion]
Inversion overnight in one subject would be detected by the subject herself, who would find that grass now looks as blood used to. So inversion is detectable after all.
\end{objection}

\begin{reply}
An overnight inversion does not preserve all relations: it breaks the relations between present experiences and memories of past ones, and the subject detects it through that break. If her memories of colors are inverted along with her experiences, she can no longer detect it. Proposition~\ref{p:undet} concerns relation-preserving inversions, which by definition leave no such trace.
\end{reply}

\section{Conclusion}

Inversion is possible because qualities are intrinsic and functional organization can be realized by physical systems with different intrinsic natures; it is undetectable because every test probes relations, and the inversion preserves them; and it leaves functionalism about the mind untouched because inverted subjects occupy the same functional roles and so have the same perceptions, beliefs and desires. The inverted spectrum is evidence about where qualitative character lies, not about what the mind is.

\begin{thebibliography}{9}\small
\bibitem[Block(1990)]{Block1990} Block, N. (1990). Inverted Earth. \emph{Philosophical Perspectives}, 4, 53--79.
\bibitem[Dennett(1988)]{Dennett1988} Dennett, D.~C. (1988). Quining qualia. In A.~J. Marcel and E.~Bisiach (eds.), \emph{Consciousness in Contemporary Science}. Oxford University Press.
\bibitem[Hardin(1988)]{Hardin1988} Hardin, C.~L. (1988). \emph{Color for Philosophers: Unweaving the Rainbow}. Hackett.
\bibitem[Locke(1690)]{Locke1690} Locke, J. (1690). \emph{An Essay Concerning Human Understanding}.
\bibitem[Palmer(1999)]{Palmer1999} Palmer, S.~E. (1999). Color, consciousness, and the isomorphism constraint. \emph{Behavioral and Brain Sciences}, 22(6), 923--943.
\bibitem[Shoemaker(1982)]{Shoemaker1982} Shoemaker, S. (1982). The inverted spectrum. \emph{The Journal of Philosophy}, 79(7), 357--381.
\end{thebibliography}

\end{document}
